\documentclass[12pt]{article}
\usepackage{amsmath, amssymb, amsthm}
\usepackage[margin=1in]{geometry}

\theoremstyle{plain}
\newtheorem{theorem}{Theorem}
\newtheorem{lemma}{Lemma}
\newtheorem{corollary}{Corollary}

\theoremstyle{definition}
\newtheorem{definition}{Definition}

\begin{document}

\title{Non-Existence of Lagrangian Smoothings for Certain Polyhedral Lagrangian Surfaces}
\date{}
\maketitle

\begin{abstract}
We examine whether a polyhedral Lagrangian surface  with the property that exactly 4 faces meet at every vertex necessarily admits a Lagrangian smoothing. We answer this question in the negative. By constructing a counterexample based on the regular octahedron, we show that such a surface can be topologically spherical. Since smooth exact Lagrangian spheres do not exist in  (a result due to Gromov), such a polyhedral surface cannot be smoothed.
\end{abstract}

\section{Introduction}

A \textbf{polyhedral Lagrangian surface}  is a finite polyhedral complex forming a topological surface, such that every 2-dimensional face is a polygon contained in a Lagrangian affine subspace of . A \textbf{Lagrangian smoothing} of  is a Hamiltonian isotopy  () of smooth Lagrangian submanifolds converging to  in the Hausdorff topology as .

The question is whether the local combinatorial condition---that the vertex degree of the 1-skeleton is 4---is sufficient to guarantee the existence of a Lagrangian smoothing.

We provide a negative answer by identifying a global topological obstruction. While the local condition allows for the construction of a polyhedral Lagrangian surface homeomorphic to the sphere , symplectic rigidity prohibits the existence of a smooth Lagrangian sphere in .

\section{Combinatorial Analysis}

We consider a polyhedral surface  satisfying the following conditions:
\begin{enumerate}
\item  is a simplicial complex (faces are triangles).
\item Exactly 4 faces meet at every vertex.
\end{enumerate}
Let  denote the number of vertices, edges, and faces of . The simplicial condition implies . The vertex degree condition implies , or .
Substituting these into the Euler characteristic formula :

For a closed connected surface, . The only integer solution for  that yields a valid Euler characteristic for an orientable surface is , which implies . This corresponds to the combinatorics of a \textbf{regular octahedron}. Thus, any simplicial polyhedral surface satisfying the vertex degree condition is homeomorphic to the sphere .

\section{Existence of the Counterexample}

To establish the answer, we must verify that a polyhedral Lagrangian octahedron actually exists in .
Let  be the vertices in . The condition that a triangular face with vertices  is Lagrangian is equivalent to the vanishing of the symplectic area: .
The octahedron has 8 faces, imposing 8 quadratic constraints on the  coordinates of the vertices. Since the system is underdetermined, solutions exist in general position (ensuring the complex is embedded and non-degenerate).
Therefore, there exists a polyhedral Lagrangian surface  homeomorphic to  satisfying the problem's conditions.

\section{Obstruction to Smoothing}

Suppose  admits a Lagrangian smoothing . By definition, for any ,  is a smooth embedded Lagrangian submanifold of . Since the isotopy is topological,  must be homeomorphic to .

We invoke the following fundamental theorem in symplectic topology:

\begin{theorem}[Gromov \cite{Gromov85}]
There are no compact exact Lagrangian submanifolds embedded in standard symplectic vector spaces .
\end{theorem}

A Lagrangian submanifold  is \textit{exact} if the restriction of the Liouville form  (where ) to  is exact. Since , any closed 1-form on  is exact. Consequently, any Lagrangian embedding of a sphere  into  is automatically exact.

The existence of a Lagrangian smoothing  would imply the existence of a smooth, exact Lagrangian sphere in , contradicting Gromov's theorem. Thus,  cannot be smoothed.

\section{Conclusion}

The polyhedral Lagrangian octahedron satisfies the condition that exactly 4 faces meet at every vertex. However, it does not admit a Lagrangian smoothing because the resulting smooth surface would violate global symplectic rigidity constraints.

\begin{thebibliography}{9}
\bibitem{Gromov85} M. Gromov, \textit{Pseudo holomorphic curves in symplectic manifolds}, Invent. Math. 82 (1985), 307--347.
\end{thebibliography}

\end{document}
